\section{Formato dei file di output}
\label{app:formato}

\subsection{\code{esperimenti.csv}}

Una riga per istanza e variante, 31 colonne, decimali con il punto. Le sequenze dei
landmark sono racchiuse fra virgolette perché contengono virgole; una
lunghezza infinita è scritta \code{INF} e la relativa sequenza \code{[]}.

\begin{table}[H]
\centering
\small
\begin{tabularx}{\textwidth}{@{}>{\raggedright\arraybackslash}p{6.1cm} X@{}}
\toprule
\textbf{Colonne} & \textbf{Significato} \\
\midrule
\code{gridId} & identificativo dell'istanza, per esempio
  \code{G12x12\_d10\_s50\_p1} (dimensioni, densità in percentuale, seed,
  numero della coppia) \\
\code{variante} & variante di \cmin{} (\cref{tab:varianti}) \\
\code{tipoGriglia} & tipo di griglia, sempre \code{Casuale} \\
\code{width}, \code{height}, \code{nOstacoli}, \code{seed} & parametri della griglia \\
\code{Ox}, \code{Oy}, \code{Dx}, \code{Dy} & coordinate effettive di $O$ e $D$ \\
\code{lunghezzaOD}, \code{lunghezzaDO} & lunghezze nei due versi, 6 decimali \\
\code{sequenzaLandmarkOD}, \code{sequenzaLandmarkDO} & sequenze dei landmark,
  ogni elemento nella forma \code{(x=..,y=..,tipo=..)} \\
\code{correttezza} & esito della verifica di simmetria \\
\code{frontieraTot} & celle di frontiera di tutte le chiusure calcolate \\
\code{falsiRiga16} & volte in cui la riga 16 è risultata falsa \\
\code{falsiRiga17} & potature dovute alla sola riga 17 \\
\code{pruningTotale} & somma delle due precedenti \\
\code{chiamateRicorsive}, \code{cacheHit} & invocazioni totali e invocazioni
  servite dalla cache \\
\code{profonditaMassima} & profondità massima della ricorsione \\
\code{chiusureAttivePicco} & matrici di chiusura attive
  contemporaneamente lungo la pila delle chiamate, al picco \\
\code{celleChiusuraTotali} & somma delle dimensioni delle chiusure calcolate \\
\code{memoriaMatriciPiccoBytesStimata} & $\code{chiusureAttivePicco} \cdot
  W \cdot H \cdot 4$ byte \\
\code{warmupEseguiti}, \code{ripetizioniTempo} & riscaldamenti e ripetizioni
  misurate \\
\code{tempoMinNs}, \code{tempoMedianoNs}, \code{tempoMaxNs} & tempi delle
  ripetizioni in nanosecondi, versi $O \to D$ e $D \to O$ insieme \\
\bottomrule
\end{tabularx}
\caption{Colonne di \code{esperimenti.csv}. Tutti i contatori sommano i due
versi.}
\end{table}

\subsection{\code{limiti\_scala.csv}}

Una riga per variante e coppia, 24 colonne.

\begin{table}[H]
\centering
\small
\begin{tabularx}{\textwidth}{@{}>{\raggedright\arraybackslash}p{6.1cm} X@{}}
\toprule
\textbf{Colonne} & \textbf{Significato} \\
\midrule
\code{variante} & \code{base} oppure \code{limite-propagato} \\
\code{width}, \code{height}, \code{densitaPerc}, \code{nOstacoli}, \code{seed}
  & parametri della griglia \\
\code{Ox}, \code{Oy}, \code{Dx}, \code{Dy} & coordinate di $O$ e $D$ \\
\code{esitoOD}, \code{esitoDO} & \code{OK}, \code{TIMEOUT}, \code{OOM},
  \code{STACK} o \code{ERRORE}; \code{esitoDO} vale \code{NON\_ESEGUITO} se
  il verso $O \to D$ non è terminato \\
\code{lunghezzaOD}, \code{lunghezzaDO} & lunghezze, \code{INF} o vuote se il
  verso non è terminato \\
\code{simmetria} & \code{true} o \code{false} se entrambi i versi sono
  terminati, altrimenti \code{ND} \\
\code{tempoODms}, \code{tempoDOms} & tempo di una singola esecuzione a
  freddo, in millisecondi \\
\code{chiamateOD}, \code{profonditaOD}, \code{frontieraOD}, \code{cacheHitOD}
  & contatori del verso $O \to D$ \\
\code{piccoHeapMBOD} & somma dei picchi delle aree di heap della JVM, in MB \\
\code{timeoutMs}, \code{heapMax} & limiti applicati a ogni verso \\
\bottomrule
\end{tabularx}
\caption{Colonne di \code{limiti\_scala.csv}.}
\end{table}

\subsection{\code{riassunti.txt}}

Una riga per ogni riga di \code{esperimenti.csv}, nel formato del riassunto
della \cref{sec:contatori}, nello stesso ordine del CSV.

\subsection{Griglie JSON}

Ogni file di \code{risultati/griglie\_json/} contiene un array di righe;
ogni elemento è lo stato della cella secondo la \cref{tab:statocella}
($0$ libera, $128$ ostacolo). Il nome del file riporta numero dello scenario,
dimensioni, ostacoli e seed, per esempio
\code{scenario\_02\_12x12\_15ost\_seed50.json}.
